% Part: many-valued-logic
% Chapter: three-valued-logics
% Section: introduction

\documentclass[../../../include/open-logic-section]{subfiles}

\begin{document}

\olfileid[pt]{mvl}{thr}{int}

\olsection{Introdução}

Se acrescentarmos apenas mais um valor~$\Undef$ a $\True$ e $\False$, obteremos uma
lógica trivalorada. Embora haja apenas um valor de verdade adicional,
as possibilidades de definir as funções de verdade de $\lnot$,
$\land$, $\lor$ e $\lif$ são bastante numerosas. Uma lógica pode então usar
qualquer combinação dessas funções de verdade, e também se pode escolher
designar apenas $\True$, ou tanto $\True$ quanto~$\Undef$.

Apresentamos aqui uma seleção das lógicas trivaloradas mais conhecidas,
suas motivações e algumas de suas propriedades.

\end{document}
